\section{Théories de Yang-Mills}

\subsubsection{Théorie de jauge en Physique:}

Soient une variété lisse $\mathcal{M}$ de dimension $m$ dont $(V, x^\mu)$ est une carte locale, un groupe de Lie $G$ et un espace de Hilbert $\mathcal{H}$ muni d'une action de $G$ sur $\mathcal{H}$ qu'on notera $U\psi$ pour $U \in G$ et $\psi \in \mathcal{H}$. Soient enfin $B_\mu : V \to \mathfrak{g}$ des applications lisse qu'on appelera champs de jauge. Et soit $g\in \R$ une constante appelée constante de couplage de la théorie. \\

Soit $U: V\to G$ une application lisse, une transformation de jauge locale du champs $\psi$ est donnée par:
$$\psi \mapsto \tilde{\psi} = U^{-1}\psi$$
On effectue alors les substitutions
$$\frac{1}{i}\frac{\partial}{\partial x^\mu} \mapsto \frac{1}{i}\frac{\partial}{\partial x^\mu} + g B_\mu$$

On impose ensuite l'invariance de jauge pour les opérateurs différentiels du premier ordre sous la transformation de jauge locale. On souhaite
$$U\left(\frac{1}{i}\frac{\partial}{\partial x^\mu} + g B'_\mu\right)\tilde{\psi} = \left(\frac{1}{i}\frac{\partial}{\partial x^\mu} + g B_\mu\right)\psi$$
On a 
\begin{eqnarray*}
U\left(\frac{1}{i}\frac{\partial}{\partial x^\mu} + g B'_\mu\right)\tilde{\psi} &=& U\left(\frac{1}{i}\frac{\partial}{\partial x^\mu} + g B'_\mu\right)U^{-1}\psi \\
&=& \frac{1}{i}U\frac{\partial U^{-1}}{\partial x^\mu}\psi + \frac{1}{i}\frac{\partial \psi}{\partial x^\mu} + gUB'_\mu U^{-1}\psi
\end{eqnarray*}
Ainsi, on doit avoir
$$\frac{1}{i}U\frac{\partial U^{-1}}{\partial x^\mu} + gUB'_\mu U^{-1} = gB_\mu$$
autrement dit
$$B'_\mu = U^{-1}B_\mu U - \frac{1}{ig}\frac{\partial U^{-1}}{\partial x^\mu}U$$
Maintenant en différentiant $U^{-1}U = I$, on a 
$$\frac{\partial U^{-1}}{\partial x^\mu}U + U^{-1}\frac{\partial U}{\partial x^\mu} = 0$$
Ce qui donne 
$$B'_\mu = U^{-1}B_\mu U + \frac{1}{ig}U^{-1}\frac{\partial U}{\partial x^\mu}$$ 

\begin{defi}
Soient $B_\mu$ et $B'_\mu$ deux champs de jauges. On dit que $B_\mu$ est $B'_\mu$ sont équivalent, s'il existe une transformation de jauge locale $U: V\to G$ telle que 
$$B'_\mu = U^{-1}B_\mu U + \frac{1}{ig}U^{-1}\frac{\partial U}{\partial x^\mu}$$
On dit que le champs $B_\mu$ est un champs de jauge pur si 
$$B_\mu = \frac{1}{ig}U^{-1}\frac{\partial U}{\partial x^\mu}$$
\end{defi}
Soit alors $F_{\mu\nu}$ le tenseur de force de la théorie donné par
$$F_{\mu\nu} = \partial_\nu B_\mu - \partial_\mu B_\nu - ig[B_\mu, B_\nu]$$
On a 
\begin{eqnarray*}
B'_\mu B'_\nu &=& \left(U^{-1}B_\mu U + \frac{1}{ig}U^{-1}\partial_\mu U\right)\left(U^{-1}B_\nu U + \frac{1}{ig}U^{-1}\partial_\nu U\right) \\
&=& U^{-1}B_\mu B_\nu U + \frac{1}{ig}U^{-1}B_\mu (\partial_\nu U) \\
&+& \frac{1}{ig}U^{-1} (\partial_\mu U) U^{-1}B_\nu U + \frac{1}{(ig)^2}U^{-1}(\partial_\mu U) U^{-1} (\partial_\nu U)
\end{eqnarray*}
Ainsi
\begin{eqnarray*}
[B'_\mu, B'_\nu] &=& U^{-1}[B_\mu, B_\nu]U \\
&+& \frac{1}{ig}U^{-1}\left( (\partial_\mu U)U^{-1}B_\nu - (\partial_\nu U)U^{-1}B_\mu \right) U \\ 
&+& \frac{1}{ig}U^{-1} B_\mu (\partial_\nu U) - \frac{1}{ig}U^{-1} B_\nu (\partial_\mu U) \\
&+& \frac{1}{(ig)^2}U^{-1}(\partial_\mu U) U^{-1} (\partial_\nu U) - \frac{1}{(ig)^2}U^{-1}(\partial_\nu U) U^{-1} (\partial_\mu U)
\end{eqnarray*}
Or 
\begin{eqnarray*}
\partial_\nu B'_\mu &=& \partial_\nu \left(U^{-1}B_\mu U + \frac{1}{ig}U^{-1}\partial_\mu U\right) \\
&=& U^{-1} (\partial_\nu U) U^{-1} B_\mu U + U^{-1} (\partial_\nu B_\mu)U + U^{-1}B_\mu (\partial_\nu U) \\
&+& \frac{1}{ig}U^{-1}(\partial_\nu U)U^{-1}(\partial_\mu U) + \frac{1}{ig} U^{-1} (\partial_\nu \partial_\mu U)
\end{eqnarray*}
Ainsi, on a
$$F'_{\mu\nu} = U^{-1} F_{\mu\nu} U$$
Qui est donc bien invariant sous une transformation de jauge locale. \\

On peut alors écrire le lagrangien de l'interaction dont on s'arrange pour qu'il soit invariant sous les transformations de jauge locale, pour obtenir des équations invariantes de jauge. Ainsi on pose
$$\mathcal{L}_{\text{libre}} = \tr \left(-\frac{1}{4}F_{\mu\nu}F^{\mu\nu}\right)$$
Finalement on peut écrire les équations d'Euler-Lagrange dans ce cas et on obtient les équations de Yang-Mills
$$\sum_\mu \partial_\mu F_{\mu\nu} + ig[B_\mu, F_{\mu\nu}] = 0$$


\subsubsection{Théorie de jauge en Géométrie:}

\begin{defi}[Théorie de jauge]
Une théorie de jauge est la donnée d'un fibré principal $G \hookrightarrow P \overset{\pi}{\to} \mathcal{M}$ et d'une connexion d'Ehresmann $\omega$. \\
Dans ce cas on appelle champs de jauge sur $P$ la connexion d'Ehresmann $\omega$ et champs de force la courbure $D\omega$. \\
Une transformation de jauge du fibré principal est un automorphisme de fibrés principaux $\psi: P \to P$ recouvrant l'identité sur $\mathcal{M}$ qui commute avec l'action à droite de $G$ sur $P$.
\end{defi}

\begin{pro}
Soit $\sigma : V\to \pi^{-1}(V)$ une section du fibré principal. On note $A = \sigma^*(\omega)$ le champs de jauge local et $F = \sigma^*(D\omega)$ le champs de force local. Alors on a
$$F = dA + \frac{1}{2}A\wedge A$$
et
\begin{alignat*}{2}
DF & = 0   && \quad\text{Identité de Bianchi} \\
D\star F & = 0  && \quad\text{Équation de Yang-Mills}
\end{alignat*}
\end{pro}

\section{Electrodynamique}

\subsubsection{Équations de Maxwell dans le langage des formes:}

Habituellement on écrit les équations de Maxwell pour le champs electromagnétique $(E, B)$ sous forme de 4 équations vectorielles:
\begin{alignat*}{2}
\Div E & =  4\pi\rho  && \quad\text{Maxwell-Gauss} \\
\rot E & = - \frac{\partial B}{\partial t}  && \quad\text{Maxwell-Farraday} \\
\Div B & =  0 && \quad\text{Maxwell-Thomson} \\
\rot B & =  4\pi j + \frac{\partial E}{\partial t}  &&\quad\text{Maxwell-Ampère}
\end{alignat*}
Et on rappelle que dans $\R^3$, on a
$$\begin{matrix}
\Omega^0 & \overset{d}{\longrightarrow} & \Omega^1 & \overset{d}{\longrightarrow} & \Omega^2 & \overset{d}{\longrightarrow} & \Omega^3\\
\Big\updownarrow &  & \Big\updownarrow &  & \Big\updownarrow\star &  & \Big\updownarrow\star\\
\R & \overset{\grad}{\longrightarrow} & \R^3 & \overset{\rot}{\longrightarrow} & \R^3 & \overset{\Div}{\longrightarrow} & \R
\end{matrix}
$$
Ainsi on remplace le champs $E = (E_x, E_y, E_z)$ par une $1$-forme 
$$\mathscr{E} = E_x dx + E_y dy + E_z dz$$ 
le champs $B = (B_x, B_y, B_z)$ par une $2$-forme 
$$\mathscr{B} = B_x dy\wedge dz + B_y dz \wedge dx + B_z dx\wedge dy$$
avec $\rho$ une $0$-forme et $j$ une $1$-forme. Les équations de Maxwell deviennent alors
\begin{alignat*}{2}
\star d(\star\mathscr{E}) & =  4\pi\rho  && \quad\text{Maxwell-Gauss} \\
d \mathscr{E} & = - \frac{\partial \mathscr{B}}{\partial t}  && \quad\text{Maxwell-Farraday} \\
d \mathscr{B} & =  0 && \quad\text{Maxwell-Thomson} \\
\star d(\star\mathscr{B}) & =  4\pi j + \frac{\partial \mathscr{E}}{\partial t}  &&\quad\text{Maxwell-Ampère}
\end{alignat*}
On se place alors dans l'espace de Minkowski $\R^{(1, 3)}$ muni de sa métrique usuelle $\eta_{\mu\nu}$ de signature $(-,+,+,+)$ et $d_{(1,3)} = d_3 + \frac{\partial}{\partial t}\wedge dt$. On pose alors $\mathscr{F} = \mathscr{E} \wedge dt + \mathscr{B}$ et $J = \rho dt + j$ la $1$-forme correspondant au quadrivecteur de courant. On a

\begin{eqnarray*}
d_{(1,3)}\mathscr{F} &=& d_{(1,3)}(\mathscr{E}\wedge dt) + d_{(1,3)}\mathscr{B} \\
&=& d_{(1,3)}\mathscr{E} \wedge dt - \mathscr{E} \wedge d_{(1,3)}dt + d_{(1,3)}\mathscr{B} \\
&=& d_3\mathscr{E} \wedge dt + d_{(1,3)}\mathscr{B} \\
&=& -\frac{\partial \mathscr{B}}{\partial t}\wedge dt + \frac{\partial \mathscr{B}}{\partial t}\wedge dt \\
&=& 0
\end{eqnarray*}
Dans $\R^{(1, 3)}$, l'étoile de Hodge donne 
$$\star(dx\wedge dy) = - dt\wedge dz \quad \star(dy\wedge dz) = - dt\wedge dx \quad \star(dz\wedge dx) = - dt\wedge dy$$
$$\star(dt\wedge dx) = dy\wedge dz \quad \star(dt\wedge dy) =  dz\wedge dx \quad \star(dt\wedge dz) = dx\wedge dy$$
$$\star(dx\wedge dy\wedge dz) = - dt \quad \star(dt\wedge dx\wedge dy) = dz \quad \star(dt\wedge dy\wedge dz) = dx \quad \star(dt\wedge dz\wedge dx) = dy$$
Ainsi on a 
\begin{eqnarray*}
\star \mathscr{F} &=& \star(\mathscr{E}\wedge dt) + \star \mathscr{B}  \\
&=& - E_x dy \wedge dz - E_y dz \wedge dx - E_z dx \wedge dy \\
& & -B_x dt \wedge dx - B_y dt \wedge dy - B_z dt \wedge dz
\end{eqnarray*}
donc
\begin{eqnarray*}
d\star \mathscr{F} &=& - d E_x \wedge dy \wedge dz - d E_y \wedge dz \wedge dx - d E_z \wedge dx \wedge dy \\
& & -d B_x\wedge dt \wedge dx - d B_y\wedge dt \wedge dy - d B_z\wedge dt \wedge dz \\
&=& -(\partial_t E_x dt + \partial_x E_x dx)\wedge dy \wedge dz \\
& & -(\partial_t E_y dt + \partial_y E_y dy)\wedge dz \wedge dx \\
& & -(\partial_t E_z dt + \partial_z E_z dz)\wedge dx \wedge dy \\
& & -(\partial_y B_x dy + \partial_z B_x dz)\wedge dt \wedge dx \\
& & -(\partial_x B_y dx + \partial_z B_y dz)\wedge dt \wedge dy \\
& & -(\partial_x B_z dx + \partial_y B_z dy)\wedge dt \wedge dz \\
&=& -(\partial_x E_x + \partial_y E_y + \partial_z E_z)dx \wedge dy \wedge dz \\
& & +(-\partial_t E_z - \partial_y B_x + \partial_x B_y) dt \wedge dx \wedge dy \\
& & +(-\partial_t E_x - \partial_z B_y + \partial_y B_z) dt \wedge dy \wedge dz \\
& & +(-\partial_t E_y - \partial_x B_z + \partial_z B_x) dt \wedge dz \wedge dx 
\end{eqnarray*}
donc 
\begin{eqnarray*}
\star d\star \mathscr{F} &=& (\partial_x E_x + \partial_y E_y + \partial_z E_z)dt \\
& & +(-\partial_t E_x - \partial_z B_y + \partial_y B_z) dx \\
& & +(-\partial_t E_y - \partial_x B_z + \partial_z B_x) dy \\
& & +(-\partial_t E_z - \partial_y B_x + \partial_x B_y) dz \\
&=& 4\pi J
\end{eqnarray*}
d'après Maxwell-Gauss et Maxwell-Ampère. \\

Ainsi les équations de Maxwell se résument en terme de formes différentielles où $\mathscr{F}$ est une $2$-forme et $J$ une $1$-forme sur $\R^{(1,3)}$ par 
$$d\mathscr{F} = 0$$
et
$$\star d\star \mathscr{F} = 4\pi J$$

